\documentclass[11pt]{article}
\usepackage[a4paper,margin=26mm]{geometry}
\usepackage{amsmath,amssymb,mathtools}
\usepackage{unicode-math}
\setmainfont{Latin Modern Roman}
\setmathfont{Latin Modern Math}
\usepackage{xurl,hyperref}
\hypersetup{hidelinks,pdftitle={Problem statement}}
\newcommand{\sref}[2]{\hyperlink{source:#1:#2}{[S#2]}}
\newcommand{\cataloguescope}[1]{#1}
\begin{document}
\section{Problem 1.20 — Justify the above arguments rigorously.}\label{statement}

\subsection{Definitions and mathematical statement}

\paragraph{Shared notation.}$\mathbb N=\{1,2,\ldots\}$, and $\mathbb Z,\mathbb Q,\mathbb R,\mathbb C$ denote the integers, rationals, reals and complex numbers. Any other conventions are specified locally.


A knot or link is a smooth oriented embedding of one or more circles in the oriented $3$-sphere, considered up to orientation-preserving ambient isotopy. Write $O$ for the unknot, which bounds an embedded disk. Write $J_N(K)$ for the Kashaev invariant $\langle K\rangle_N$, equivalently the reduced $N$-dimensional colored Jones polynomial evaluated at $q=\exp(2\pi i/N)$, with $J_N(O)=1$. Reduction removes the unknot quantum-dimension factor before evaluation at the root of unity. For $q=\exp(2\pi i/N)$ put $(q)_n=\prod_{j=1}^n(1-q^j)$ and $(q)_0=1$. For the knot $5_2$ the source starts from
\[J_N(5_2)=\sum_{0\le m\le n<N}\frac{(q)_n^2}{\overline{(q)_m}}q^{-m(n+1)}.\]
Let $\operatorname{Li}_2(z)=-\int_0^z\log(1-s)\,ds/s$, with branches chosen consistently along the domains under study. The formal continuum potential is
\[\Phi(z,w)=\frac{\pi^2}{2}-2\operatorname{Li}_2(z)-\operatorname{Li}_2(1/w)-\log z\log w,\]
initially on $|z|=|w|=1$, $0\le\arg w\le\arg z\le2\pi$.
\paragraph{Statement recorded in the supplied source.}
Make rigorous the source's proposed asymptotic passage from the finite root-of-unity sum to an oscillatory double integral and then to a dominant critical value of $\Phi$. This requires justified product asymptotics, branches and phases, error bounds sufficient at exponential scale, control of cancellations and endpoints, and any necessary contour deformation. Establish the resulting volume asymptotics for $5_2$, and justify the corresponding critical-value method for the analogous state sums of general hyperbolic knots described by Yokota. A formal equality of a critical value with volume by itself does not control the original discrete sum. \sref{statement}{2}
\subsection{Short English statement}
Justify rigorously the product, sum-to-integral and saddle-point steps proposed for root-of-unity colored Jones asymptotics, including the $5_2$ example and the general hyperbolic-knot method.
\subsection{Sources}
\begin{description}
\item[\hypertarget{source:statement:1}{S1}] ULAM. AI and the UnsolvedMath Contributors, \emph{UnsolvedMath: A Curated Collection of Open Mathematics Problems}, catalogue entry. CC BY 4.0. \url{https://huggingface.co/datasets/ulamai/UnsolvedMath}.
\item[\hypertarget{source:statement:2}{S2}] T. Ohtsuki (editor), \emph{Problems on invariants of knots and 3-manifolds}, Geometry \& Topology Monographs 4 (2002), pp. 377--572 (published 2004), item 1.20 and its surrounding definitions and remarks. \url{https://arxiv.org/abs/math/0406190v1}.
\end{description}
\label{end:statement}
\end{document}
